% ---------------------------------------------------------------------------
% howtoread.tex
% ---------------------------------------------------------------------------
\chapter*{How to Read This Booklet}
\addcontentsline{toc}{chapter}{How to Read This Booklet}
\markboth{How to Read This Booklet}{How to Read This Booklet}

\section*{Typographic conventions}

\begin{itemize}
  \item \code{Monospaced text} is code: an identifier, a keyword, a file name
        or a command.
  \item \kw{Coloured monospaced text} marks a language keyword under
        discussion.
  \item \file{Green monospaced text} is a file in the accompanying source
        tree.
  \item \tool{Sans-serif} names are tools: \tool{QuestaSim}, \tool{Verilator},
        \tool{cocotb}, \tool{GTKWave}.
\end{itemize}

\section*{The three kinds of box}

\begin{gotcha}{What a trap looks like}
A trap marks a mistake you are statistically likely to make.
\end{gotcha}

\begin{tipbox}{What a recipe looks like}
A recipe is a pattern to copy. When the text says ``write it like this'',
this is the box it is in.
\end{tipbox}

\begin{notebox}{What a note looks like}
A note supplies background: where a construct came from, what the standard
actually says, why a tool behaves the way it does. Skip these on a first
reading and come back when something surprises you.
\end{notebox}

\begin{ruleofthumb}
A rule of thumb is one line you should be able to recite from memory.
\end{ruleofthumb}

\section*{If your simulator is Verilator, or you have no licence at all}

Read \cref{ch:verilator} from the start; \cref{app:tooling} lists exactly
what to install first. Everything in it works with free tools on your own
machine.

\section*{If your golden model is already written in Python}

Go straight to \cref{ch:python} and use \tool{cocotb}. Every technique in
\refverificationbook's layered testbench works in any IEEE 1800 simulator,
including \tool{Verilator}; this booklet only adds Python as the language
driving it.

\section*{A word on tools}

This booklet assumes you have access to a commercial simulator, most likely
\tool{QuestaSim}, because that is what a university lab has and what an
employer will hand you. It also assumes you would rather not depend on one.
\Cref{app:tooling} lists what to install and how; \cref{app:simcommands}
gathers every command line used in the booklet onto one page.

\section*{The companion booklets}

This is the \textbf{Simulation Cookbook}. \refdesignbook and
\refverificationbook cover the language this booklet's testbenches are
written in. \refexamplebook runs every tool discussed here against one
complete design.
